\chapter{The methods, one by one}\label{ch:methods}

Each section follows the same pattern: the \emph{question}, the \emph{idea}, a \emph{worked toy example} you can check by hand, \emph{where it lives in the code}, and \emph{pitfalls}. The toy examples are real scripts in \file{notes/guide/book/snippets/}; their printed output is quoted.

\section{H1: comment level, Mantel--Haenszel}\label{sec:h1}

\textbf{Question.} On a comment that was replied to, do the repliers downvote it more than other voters on that same comment?

\textbf{Idea.} For each comment, make a $2\times2$ table:
\begin{center}
\begin{tabular}{@{}l c c@{}}
\toprule
 & downvoted & upvoted \\
\midrule
voter \emph{replied} to the comment & $a$ & $b$ \\
voter \emph{did not} reply & $c$ & $d$ \\
\bottomrule
\end{tabular}
\end{center}
Within one comment the odds ratio is $ad/(bc)$. Comments differ hugely in size and tone, so we cannot just add all tables together; Mantel and Haenszel proposed pooling the tables with weights so that each comment contributes in proportion to its information:
\[
\widehat{\mathrm{OR}}_{MH}=\frac{\sum_k a_kd_k/n_k}{\sum_k b_kc_k/n_k},\qquad n_k=a_k+b_k+c_k+d_k .
\]
Intuition: the numerator adds up the ``replier downvoted and other upvoted'' evidence, the denominator the opposite, each scaled by the table size; comments with few votes automatically weigh little. \src{Mantel \& Haenszel 1959}.

\textbf{Toy example.} Two comments, X and Y.
\begin{center}
\begin{tabular}{@{}l cccc c@{}}
\toprule
 & $a$ & $b$ & $c$ & $d$ & OR inside \\
\midrule
X & 6 & 2 & 24 & 16 & 2.0 \\
Y & 1 & 9 & 4 & 36 & 1.0 \\
\bottomrule
\end{tabular}
\end{center}
Stacking the tables gives $A=7,B=11,C=28,D=52$ and a crude OR of $7\cdot52/(11\cdot28)=1.18$. The Mantel--Haenszel OR is $(2+0.72)/(1+0.72)=1.58$. The script \file{snippets/mh\_toy.py} prints exactly this. The crude value lies between the two within-comment ratios (2 and 1) but is further from 2 than the MH value: stacking mixes a heavily downvoted comment (X) with a mostly upvoted one (Y), and that mix distorts the comparison. Choosing the stratified estimate is a choice to compare like with like. With real data the gap is larger, and can even change the sign (Simpson's paradox).

\lstinputlisting{snippets/mh_toy.py}

\textbf{Uncertainty.} The interval uses the variance estimator of Robins, Breslow and Greenland, which stays valid when most tables are tiny (as ours are, a few votes per comment) \src{Robins et al.\ 1986}. In addition we bootstrap over comments with Poisson weights.

\textbf{In the code.} \file{04\_h1\_comment\_level.py}: a DuckDB query builds, per comment, the four counts $a,b,c,d$; \texttt{mh()} computes the estimate and the variance; a loop applies the Poisson weights. Yearly results are the points in Figure~\ref{fig:h1}. A copy of the function is in \file{08\_positive\_votes.py} for upvotes (columns swapped).

\textbf{Result.} $\widehat{\mathrm{OR}}_{MH}=4.43$ [4.42, 4.44]; repliers' downvote share 53.3\% against 30.1\%; the OR grows from 3.5 for top-level comments to 9.0 at greater depth.

\begin{pitfall}
\textbf{Different voters.} Repliers and non-repliers on a comment are \emph{different people}. If repliers are, in general, heavier downvoters, part of the 4.43 is a trait of repliers, not of the act of replying. A diagnostic that uses each voter as their own control (voter as stratum) gives OR $=3.49$, and 89\% of voters with at least 10 votes on both sides downvote more on comments they reply to. So the effect survives, reduced. The confidence interval in the table above also treats comments as independent although many votes come from the same voters; the thesis should report it as ``too narrow'' (limitation M-1).
\end{pitfall}

\begin{tryit}
Run \file{mh\_toy.py}. Change table Y so that its inside OR is 2 as well. What should the MH estimate be? Then change the sizes of X and Y and see how the weights move the estimate.
\end{tryit}

\section{The null model: who is the ``random alternative''?}\label{sec:null}

\textbf{Question.} In H2 we want to know whether $A$ has a history of voting against $B$. But $A$ votes against many people. Is the figure for $B$, the person $A$ answered, unusually negative?

\textbf{Idea.} Compare with a person $A$ \emph{could have} replied to but did not. For each real reply by $A$ in article $S$, draw as the null target the author of a \emph{uniformly random earlier comment in the same article}. This keeps the article (topic, size, period) and the exposure (that person was present and earlier) and changes only the choice. The draw is deterministic: it is a hash of the comment identifier and a fixed seed (20261006), so reruns give the same null.

\textbf{In the code.} \file{02\_null\_pairs.py} produces one null target per real edge. Self-targets and deleted authors are excluded as for real edges.

\begin{pitfall}
\textbf{Article-level, not thread-level.} Many replies are in the same thread; a person who replies to a thread may be more likely to meet its authors than a random article participant. The article-level null is therefore \emph{easy to beat}. You decided to keep it and add a thread-level null next to it. The second one is designed (Chapter~\ref{ch:critique}) but not run.
\end{pitfall}

\textbf{A stronger null from the group's own models.} \src{G\'omez et al.\ 2013} model how a thread grows: the next comment attaches to an earlier one with a probability that depends on its \emph{popularity} (how many replies it already has), its \emph{recency} (attractiveness decays with the number of comments posted since) and a \emph{root bias} (extra weight on the original post). Time is measured in comment arrivals, not clock time. Lan Ceqin's thesis fits an extension that adds a \emph{reciprocity} term (replying back to someone who replied to you) separately for each period, on a sample of about 684 threads per period. If we fitted such a model on DerStandard threads and drew the null target from it, the null would respect popularity, recency and the root, which our uniform draw ignores. It would be a harder baseline to beat. Fitting needs only hundreds of threads per period, so it is cheap; it is a candidate for the next round of analyses, not done yet. Keep the three nulls side by side (article-level uniform, thread-level uniform, model-based) and report how the +13.3 points change.

\section{H2: dyads and the Beta posterior}\label{sec:h2}

\textbf{Question.} Is $A$ in general for or against $B$, and are repliers more often against the person they answer than against a null target?

\textbf{Idea in two steps.}
\begin{enumerate}
  \item \emph{Pair relation.} For the pair $(A,B)$ count how often $A$ voted up ($u$) and down ($v$) on $B$'s comments, \emph{excluding} the replied comment itself in the ``excl'' variant (so the reply does not explain its own vote). Treat ``A votes positively on B'' as a coin with unknown probability $p$. Start from a Beta$(\alpha,\beta)=(1,2)$ prior and update with the counts. The posterior mean is $E[p]=(u+\alpha)/(u+v+\alpha+\beta)$.
  \item \emph{Edge sign.} If $E[p]<0.4$ the relation is \emph{negative}, if $>0.6$ \emph{positive}, otherwise no edge. The headline restricts to pairs with at least 5 votes.
\end{enumerate}
The prior and thresholds are inherited from FAULTANA so that our relations are comparable with theirs \src{Fraxanet et al.\ 2024}. The model is the standard beta--binomial one \src{Gelman et al.\ 2013}.

\textbf{Toy example.} \file{snippets/beta\_shrink\_toy.py} prints how the posterior mean behaves:
\begin{center}
\begin{tabular}{@{}rrcc@{}}
\toprule
$u$ & $v$ & $E[p]$ & edge \\
\midrule
1 & 3 & 0.286 & negative \\
2 & 8 & 0.231 & negative \\
50 & 50 & 0.495 & none \\
400 & 100 & 0.797 & positive \\
0 & 1 & 0.250 & negative (from one vote!) \\
1 & 0 & 0.500 & none \\
\bottomrule
\end{tabular}
\end{center}
The last two rows show why a minimum number of votes is needed: a single downvote already yields a ``negative'' edge, because the prior leans negative. With a lean-negative prior the rule is slightly more generous to negative edges; the multiverse (Section~\ref{sec:multiverse}) checks other priors.

\textbf{Results (Figure~\ref{fig:h2}).} Among defined pairs, 31.4\% of real reply pairs are negative against 18.1\% of null pairs, a difference of $+13.3$ percentage points with a 95\% interval of [12.6, 14.0]. $A$'s share of downvotes among all votes to the target is 6.9\% (real) against 4.5\% (null).

\textbf{Uncertainty.} Paired Poisson bootstrap over repliers, 200 replicates: each replier keeps their real and null edges together and gets one weight.

\textbf{In the code.} \file{03\_dyads.py} aggregates votes by pair; \file{05\_h2\_h3\_dyad\_faction.py} computes the Beta posterior, the edge sign, the comparison against the null and the bootstrap.

\begin{pitfall}
\textbf{The defined-edge subsets differ.} The share of negative edges is computed among pairs with $\ge5$ votes. Real pairs and null pairs do not meet this criterion equally often (repliers have more history with the people they answer), so the 31.4\% and 18.1\% are computed on different subsets. The share of downvotes among \emph{all} votes to the target (6.9 vs 4.5\%) avoids the problem. A design with equal coverage is on the list of fixes (M-2).
\end{pitfall}

\section{H3: factions}\label{sec:h3}

\textbf{Question.} FAULTANA assigns 14,817 long-lived users to one of two factions. Do replies cross factions less than random mixing, and are cross-faction replies more negative?

\textbf{Idea.} If a share $p$ of users is in one faction, random mixing would produce cross-faction pairs with probability $2p(1-p)$. Compare (i) the observed share of cross-faction replies, (ii) the share among null pairs and (iii) random mixing. Then, among replies with both users in the data, compare how often the replier downvoted the parent within a faction and across factions.

\textbf{Results.} Cross-faction share: 26.0\% of real replies, 23.2\% of null pairs, 46.8\% under random mixing. So replies are strongly \emph{within} faction relative to random mixing, but slightly \emph{more cross-faction than the null}. When the replier is in the other faction they downvoted the parent in 64.6\% of cases against 37.0\% within a faction.

\begin{pitfall}
\textbf{Partial circularity.} The factions were derived from the same votes we analyse, so ``cross-faction pairs have more negative votes'' is partly true by construction. Treat the faction results as a consistency check, not independent evidence. The group's own projection of factions to all users (AUC 0.63) was too weak to use.

\textbf{Do the two factions beat a proper baseline?} \src{Leskovec et al.\ 2010} looked for a two-faction split in three signed networks and compared the share of edges the best partition satisfies (positive inside, negative between) with two baselines: the same network with \emph{shuffled signs}, and a random network that keeps each node's number of positive and negative edges. They found the real partitions satisfied no more edges than the baselines, so no evidence of global two-camp structure. FAULTANA's partition satisfies 71\% of edges (29\% frustrated) while a single camp would satisfy only 59\% (41\% of edges are negative), so it is well above the trivial solution, but we did not check it against shuffled or degree-preserving baselines. A committee member who knows that paper could ask. Also, \src{Easley \& Kleinberg 2010} show that \emph{weak} balance allows several groups, not two, so the two-faction assumption is itself a choice.
\end{pitfall}

\section{H4: lifetime}\label{sec:h4}

\textbf{Step 1: define lifetime.} Lifetime $=$ last activity $-$ first activity, where activity is a comment \emph{or} a vote. A user is right-censored if the last activity is after 2022-07-01, left-censored if the first activity is before 2013-07-01. 26.7\% of users are active for less than a day, a one-time visitor mass that dominates the left part of Figure~\ref{fig:km}.

\textbf{Step 2: Kaplan--Meier by cohort.} Entry year is the cohort. The curves in Figure~\ref{fig:km} fall quickly in the first weeks and then flatten; median lifetimes are 1.36, 0.81, 0.41 and 0.39 years for cohorts 2014, 2016, 2018 and 2020. Newer cohorts leave faster.

\textbf{Toy example.} \file{snippets/km\_toy.py}: six users, times 1, 2, 2, 4, 5, 6 years, the second 2, the fifth and the sixth still active (censored). Output:
\begin{center}
\begin{tabular}{@{}cccc@{}}
\toprule
$t$ & at risk & leave & $S(t)$ \\
\midrule
1 & 6 & 1 & 0.833 \\
2 & 5 & 1 & 0.667 \\
4 & 3 & 1 & 0.444 \\
\bottomrule
\end{tabular}
\end{center}
At $t=2$ the user censored at 2 is still counted at risk (5 people), but is not counted as leaving; after $t=2$ that person disappears from the risk set. At $t=4$ only 3 remain at risk.

\lstinputlisting{snippets/km_toy.py}

\textbf{Step 3: the landmark cohort.} For H4 we keep users still active at day 30 (98,670 users), measure their first 30 days (downvote share received, hostile replies received and written, whether they got any reply, number of comments on the log scale) and follow them from day 30.

\textbf{Step 4: shrinkage.} A user with one reply has a ``hostile share'' of 0\% or 100\%, which is noise. We pull each share toward the population share $\bar p$ with a pseudo-count $M=2$: $(k+M\bar p)/(n+M)$. With $\bar p=0.17$, \file{snippets/beta\_shrink\_toy.py} gives:
\begin{center}
\begin{tabular}{@{}lc@{}}
\toprule
Hostile $k$ of $n$ & shrunk share \\
\midrule
0 of 0 & 0.170 \\
1 of 1 & 0.447 \\
0 of 1 & 0.113 \\
5 of 10 & 0.445 \\
50 of 100 & 0.494 \\
\bottomrule
\end{tabular}
\end{center}
So one hostile reply out of one is treated like 5 of 10, but 50 of 100 is almost taken at face value. This is an empirical-Bayes idea \src{Gelman et al.\ 2013}. An earlier bug here (zero replies giving an undefined share) is choice C14.

\textbf{Step 5: Cox model.} Features are standardised (hazard ratio per standard deviation), baselines are stratified by cohort, and a small ridge penalty (0.001) stabilises estimation. Results (Figure~\ref{fig:cox}): downvote share received HR 1.11; hostile replies received 1.06; hostile replies written 1.04; having received any reply 0.95; log number of early comments 1.23.

\textbf{What the literature adds.} Three points from the second reading pass (Section F of the papers summary). (i) \src{Danescu-Niculescu-Mizil et al.\ 2013} found that a plain activity rate (time between posts) and the calendar month predict who leaves better than their language features; our log early comments (HR 1.23) plays that role. The informative figure is therefore the hazard ratio of the vote-based covariates \emph{with and without} the activity covariate in the model, to show what they add. (ii) \src{Cheng et al.\ 2015} found that the early signals that best predicted a later ban were a low proportion of up-votes received and a high posting rate per day; they also used the largest number of posts in a single thread, a burst measure we do not have. (iii) \src{Wang et al.\ 2022} remind us that the landmark method analyses only users who survive to the landmark, so H4 describes the 98,670 users still active at day 30, not all commenters.

\begin{pitfall}
\textbf{Read hazard ratios correctly.} 1.23 for early comments means that people who comment more in the first 30 days leave faster \emph{per unit time, after day 30}. It does not mean that commenting makes people leave; heavy early commenters include intense, short-burst users. And the votes behind the downvote share are not timed: some may have been cast after day 30 (limitation H-3).
\end{pitfall}

\section{Multiverse: how fragile are the headlines?}\label{sec:multiverse}
\begin{idea}
Re-run the analysis under every defensible combination of arbitrary choices and look at the spread; a result that flips sign under a reasonable choice is not robust.
\end{idea}
\src{Steegen et al.\ 2016} call this a multiverse analysis; \src{Simonsohn et al.\ 2020} formalise it as specification curves and add a joint test. We do the descriptive part only: 24 specifications for H2 (negative-edge gap between $+11.7$ and $+20.6$ percentage points, Figure~\ref{fig:multiverse}) and 18 for H4 (hazard ratios between about 1.10 and 1.12 for the downvote share and 1.05 to 1.08 for hostile replies). All specifications keep their sign, so the headline direction is not an artefact of one threshold. Simonsohn's joint inference has not been done, so the thesis may say ``consistent across specifications'' but not ``jointly significant''.

\textbf{How the joint test works.} The null hypothesis is that the effect is zero in \emph{every} specification. There is no formula for it because the specifications share the same data, so it is obtained by resampling. For observational data: (1) estimate all $K$ specifications and keep their estimates $\hat b_k$; (2) remove each estimated effect from its own outcome, $y^*_k=y_k-\hat b_k x_k$, so the null is true by construction; (3) draw $N$ rows with replacement, \emph{using the same drawn rows for all $K$ specifications}; (4) re-estimate all $K$; (5) repeat 500--1000 times; (6) compare the observed statistic with its resampled distribution. Three statistics are used: the median estimate, the share of specifications that are significant in the predicted direction, and the average $Z$ (Stouffer). For H2 the natural resampling unit is the replier, as in our current bootstrap; for H4 the user. Expect tens of minutes of compute. \src{Simonsohn et al.\ 2020}, Supplement 4, shows that this also repairs specifications with inflated false-positive rates, and that the combined test has more power than any single specification.

\section{Positive votes: the mirror analysis}\label{sec:positive}
You asked whether the same logic works for upvotes. \file{08\_positive\_votes.py} repeats H1--H4 with the roles swapped. By construction the H1 result is the inverse of the downvote one ($0.23=1/4.43$). New information comes only from H2--H4. Repliers upvote the target $93.1\%$ of the time against $95.5\%$ for null targets; positive edges are 68.6\% against 81.9\% (a mirror of $-13.3$ pp); upvote share received has HR 0.87 (more upvotes, lower hazard of leaving); friendly replies have HR about 1.00 (nothing).

\section{Long-lived users}\label{sec:longlived}
\textbf{Question.} Do the top 10\% longest-living users behave differently in their first month?

\textbf{Result.} In the first 30 days the top 10\% within cohort receive fewer downvotes (20.4\% against 30.0\%), receive fewer hostile replies, and write fewer hostile replies (12.2\% against 17.1\%). In a logistic regression with cohort fixed effects, the odds ratio of being in the group is 0.87 per standard deviation of the early profile.

\begin{pitfall}
\textbf{This is a selected group.} Lifetime is last minus first activity, so the top 10\% of a cohort is mostly made of users who are \emph{still active} when the data end (right-censored survivors), ordered by when they were last seen. ``Long-lived'' mixes users who truly stay with users who joined early enough to be seen for long. The profile differences are real, but the group definition should be improved before the thesis states a conclusion (limitation H-2). Three options were discussed:
\begin{description}
  \item[A] top 10\% within each cohort (current);
  \item[B] ``survived $\ge3$ years'' among cohorts that can be observed that long (recommended as the headline);
  \item[C] top 10\% among early cohorts only, keeping the wording.
\end{description}
\end{pitfall}

\section{Reply-back and tenure tables}
Two descriptive tables accompany H2. \emph{Reply-back:} when $A$ downvoted $B$'s comment and $B$ answers, how often does $B$ reply back? (31.8\% if $A$ had downvoted, 29.8\% if no vote, 16.6\% if upvoted.) \emph{Tenure:} veterans answering newcomers downvote the parent in 68\% of cases against 48\% for veterans answering veterans. Neither has an uncertainty estimate beyond the large sample.

\section{Summary table}\label{sec:headline}
\begin{tabularx}{\textwidth}{@{}l Y l@{}}
\toprule
 & Headline & Main caveat \\
\midrule
H1 & Repliers downvote their parent 53\% vs 30\% (OR 4.43; 3.49 when each voter is own control) & Different voters; CI too narrow \\
H2 & Negative pair relations 31.4\% vs 18.1\% (+13.3 pp); consistent across 24 specs & Article-level null; unequal coverage \\
H3 & Replies within faction vs random mixing; cross-faction replies more negative & Partly circular \\
H4 & Early downvote exposure HR 1.11 per SD & Untimed votes; observational only \\
Long-lived & Top 10\% received fewer downvotes in month 1 & Censored survivors group \\
\bottomrule
\end{tabularx}
